\section{Proof Sketch}\label{sec:proof-sketch}

The proof first converts the polynomial coefficient list of $B$ into a uniform operator bound and then
differentiates $F/\|F\|_2$; this yields the static and homogeneous projective certificate in
Lemmas~\ref{lem:step-001-height} and Proposition~\ref{prop:step-001-projective}. Persistent affine
combinations form a proper affine locus and are null under every full-joint density by
Proposition~\ref{prop:step-002-law-null}.

For a measurable pivot partition, the coefficient sweep is exhausted by regions on which the pivot is at
least $1/n$. The equal-dimensional area formula, with multiplicity, controls every finite chart and monotone
convergence removes $n$; Proposition~\ref{prop:step-003-pivot-sweep} gives both chart inequalities. A
fixed-section change of variables then cancels the pivot denominator exactly and produces the normal-velocity
integral in Proposition~\ref{prop:step-004-s2-affine-swept-area}.

Brunn--Minkowski centralizes each translated cube section, and Ball's theorem controls its central section by
$\sqrt2(2R)^{N-1}$ \citep{ball1986cube}. The derivative closure controls the affine normal velocity. Their
product gives the coordinate-free affine rate and ordered capacity supremum in
Proposition~\ref{prop:step-007-s2-general-affine-rate}. In the homogeneous case the radial derivative vanishes
on the actual root section, leaving only projective speed; Proposition~\ref{prop:step-008-s2-pairwise-homogeneous-rate}
and Proposition~\ref{prop:step-008-s2-pf-closure} give the sharper coefficient.

The monic specialization is handled inside the same affine sweep. The constant and outer monomial pivots give
the two location-free velocity estimates, while the leading coefficient remains deterministic. The complete
presentation, certificate, and probability recovery appear in
Proposition~\ref{prop:step-010-complete-affine-monic-wrapper}. Finally, the one-entry shear is differentiated
directly, and both coefficient wedges are integrated exactly; Proposition~\ref{prop:step-011-scale-comparison}
keeps the lower ratio, all-law upper coefficient, and geometric section normalization distinct. The final
conjunction is Proposition~\ref{prop:step-012-anchored-coefficient-sweep}.
